\documentclass[11pt,a4paper]{article}
\usepackage[margin=2.2cm]{geometry}
\usepackage{ctex}
\usepackage{amsmath,amssymb,amsthm}
\usepackage{enumitem}
\usepackage{tasks}
\usepackage{array,tabularx}
\usepackage{multicol}
\usepackage{fancyhdr}
\usepackage{tikz}
\usepackage{graphicx}
\usepackage{xcolor}
\usepackage{needspace}

\pagestyle{fancy}
\fancyhf{}
\fancyhead[L]{\small 大模型多模态 期末试卷}
\fancyhead[R]{\small 第 \thepage\ 页}
\renewcommand{\headrulewidth}{0.4pt}

\setlength{\parindent}{0pt}
\setlength{\parskip}{0pt}

% 工具命令
\newcommand{\blank}[1]{\underline{\hspace{#1}}}
\newcommand{\fn}[2]{\dfrac{#1}{#2}}

% 大题标题
\newcommand{\examsection}[2]{%
  \vspace{0.6cm}
  \noindent{\CJKfamily{SimHei}\bfseries\large #1\quad (#2)}\par
  \vspace{0.3cm}
}

% 题号 + 题干 一行写完（避免 parskip 陷阱）
\newcommand{\qitem}[2]{%
  \noindent\textbf{#1.}\;#2\par
  \vspace{0.3em}
}

% 选择题选项（4个一行/两行可控）
\newcommand{\fourchoices}[4]{%
  \begin{enumerate}[label=(\Alph*),leftmargin=2.2em,itemsep=0pt,parsep=0pt,topsep=0pt]
    \setlength{\itemindent}{0pt}
    \item #1
    \item #2
    \item #3
    \item #4
  \end{enumerate}
  \vspace{0.4em}
}

% 选择题选项（AB 一行，CD 一行；各占 1/2 宽度，区域内左对齐）
\newcommand{\fourchoicestwo}[4]{%
  \noindent
  \begin{minipage}[t]{0.5\linewidth}\makebox[2.2em][l]{(A)}#1\end{minipage}%
  \begin{minipage}[t]{0.5\linewidth}\makebox[2.2em][l]{(B)}#2\end{minipage}\par
  \noindent
  \begin{minipage}[t]{0.5\linewidth}\makebox[2.2em][l]{(C)}#3\end{minipage}%
  \begin{minipage}[t]{0.5\linewidth}\makebox[2.2em][l]{(D)}#4\end{minipage}\par
  \vspace{0.4em}
}

% 两栏式选项（短选项）
\newcommand{\twochoices}[4]{%
  \begin{tasks}[counter-format={(A)},label-align=left,label-offset={0.4em},label-width={1.6em},item-indent=2em,column-sep=1.5em](2)
    \task #1
    \task #2
    \task #3
    \task #4
  \end{tasks}
  \vspace{0.4em}
}

% 答题区横线
\newcommand{\answerlines}[1]{%
  \vspace{0.2cm}
  \foreach \i in {1,...,#1} {
    \noindent{\color{white}\rule{\linewidth}{0.4pt}}\par\vspace{0.5cm}
  }
}

\begin{document}

% ============ 试卷标题与说明 ============
\begin{center}
{\zihao{2}\CJKfamily{SimHei}\bfseries 大模型多模态 期末试卷}\\[0.3cm]
{\zihao{4}\kaishu 满分 150 分 \quad 考试时间 150 分钟}
\end{center}

\vspace{1.5cm}
\noindent\rule{\linewidth}{1pt}

\vspace{0.2cm}
\noindent{\CJKfamily{SimHei}\bfseries 注意事项：}
\begin{enumerate}[leftmargin=2em,itemsep=0pt,parsep=0pt,topsep=0pt]
  \item 本试卷共 \textbf{三} 大题，包括选择题、填空题、解答题，共 \textbf{24} 小题。
  \item 请将答案写在答题区指定位置；解答题须写出必要的文字说明、推导过程或演算步骤。
  \item 涉及数学公式推导时，请清晰写出关键中间步骤；只写最终结果不得分。
  \item 本试卷中如无特殊说明，图像分辨率为 $H\times W$，patch 大小为 $P$，视觉 token 数为 $N_v$，文本序列长度为 $L$，隐藏维度为 $d$，对比学习温度系数为 $\tau$，扩散步数为 $T$，批大小为 $B$。
\end{enumerate}

\vspace{0.2cm}
\noindent\rule{\linewidth}{0.6pt}

% ============ 一、选择题 ============
\examsection{一、选择题}{本大题共 12 小题，每小题 5 分，共 60 分。在每小题给出的四个选项中，只有一项是符合题目要求的。}

\qitem{1}{关于多模态大模型的三个核心问题——表示（Representation）、对齐（Alignment）与融合（Fusion），下列说法\textbf{正确}的是}{}
\fourchoices{表示关注如何将不同模态映射到统一语义空间；对齐关注跨模态语义对应关系；融合关注多模态信息在下游任务中的联合利用}{表示、对齐、融合三者完全等价，只是不同论文的命名差异}{融合仅发生在预训练阶段，推理时不再需要跨模态交互}{对齐问题仅存在于图文检索任务，在 VQA 等生成任务中不存在}

\qitem{2}{关于 CLIP（Contrastive Language-Image Pre-training）的核心训练机制，下列说法\textbf{正确}的是}{}
\fourchoices{CLIP 采用单塔编码器，将图像与文本拼接后做自注意力融合}{CLIP 使用双塔结构分别编码图像与文本，通过批内对比学习（InfoNCE）拉近匹配图文对、推远不匹配对}{CLIP 的训练目标是最小化图像与文本在像素空间的均方误差（MSE）}{CLIP 只能处理固定 224$\times$224 图像，无法泛化到其他分辨率}

\qitem{3}{设 Vision Transformer（ViT）输入图像分辨率为 $224\times 224$，patch 大小为 $16\times 16$，并额外添加 1 个 CLS token。则送入 Transformer 编码器的序列长度为}{}
\fourchoicestwo{$14$}{$196$}{$197$}{$225$}

\qitem{4}{下列多模态大模型架构与其核心连接器（Connector）的匹配，\textbf{正确}的是}{}
\fourchoices{LLaVA — Q-Former；BLIP-2 — 线性投影层；Flamingo — Perceiver Resampler}{LLaVA — 线性/MLP 投影层；BLIP-2 — Q-Former；Flamingo — 门控交叉注意力层（Gated Cross-Attention）}{三者均采用 Q-Former 作为唯一连接器}{LLaVA — 扩散 U-Net；BLIP-2 — VAE 编码器；Flamingo — CLIP 文本塔}

\qitem{5}{关于 Early Fusion、Late Fusion 与 Intermediate Fusion（交叉注意力融合），下列判断\textbf{错误}的是}{}
\fourchoices{Early Fusion 在输入层即将多模态特征拼接或相加，计算简单但模态异质性处理困难}{Late Fusion 各模态独立编码后在决策层融合，模块化强但难以捕获细粒度跨模态交互}{Intermediate Fusion 通过 Cross-Attention 让一种模态的 Query 查询另一种模态的 Key/Value，是 LLaVA、Flamingo 等的主流方案}{Late Fusion 一定优于 Intermediate Fusion，因为避免了跨模态注意力带来的 $O(N_v^2)$ 计算开销}

\qitem{6}{Stable Diffusion 属于 Latent Diffusion Model（LDM），其核心设计是}{}
\fourchoices{直接在原始像素空间 $x\in\mathbb{R}^{H\times W\times 3}$ 上进行 $T$ 步去噪}{先用 VAE 将图像压缩到低维隐空间 $z$，再在隐空间训练扩散去噪网络，大幅降低计算量}{完全放弃 VAE，仅使用 GAN 判别器指导扩散过程}{只在文本嵌入空间做扩散，不生成任何图像}

\qitem{7}{关于 DDPM 前向加噪过程 $q(x_t|x_{t-1})=\mathcal{N}(\sqrt{1-\beta_t}\,x_{t-1},\,\beta_t I)$，下列说法\textbf{正确}的是}{}
\fourchoices{前向过程是非马尔可夫的，当前状态依赖全部历史步}{当 $t\to T$ 且 $\bar\alpha_t=\prod_{s=1}^{t}(1-\beta_s)\to 0$ 时，$x_T$ 近似服从标准高斯分布}{反向去噪过程 $q(x_{t-1}|x_t)$ 有解析闭式解，可直接采样}{方差调度 $\beta_t$ 必须恒定为常数，否则训练不稳定}

\qitem{8}{关于 VQ-VAE / VQ-GAN 中离散视觉 Token 化的机制，下列说法\textbf{正确}的是}{}
\fourchoices{Codebook 中每个向量通过 argmax 操作被选中，梯度可通过 Straight-Through Estimator 近似回传}{离散化后图像无法用 Transformer 建模，必须退回 CNN}{VQ 损失仅包含重建损失，不含 codebook 与 commitment 项}{Codebook 大小对生成质量无影响，越大越好且无需担心码本利用率}

\qitem{9}{关于 Flux 文生图模型的底层技术与特性，下列说法\textbf{错误}的是（\quad）}{}
\fourchoices{Flux 基于 Flow Matching 流匹配范式训练，而非传统 DDPM 扩散损失}{Flux 主干网络采用 DiT（Diffusion Transformer），抛弃了 Stable Diffusion 经典 U-Net 结构}{Flux 原生对图像内文字、人体手部细节生成效果显著优于 SDXL}{Flux 基础版本显存占用低于 SDXL，8G 显存显卡可无压力本地运行}

\qitem{10}{在 Whisper 音频编码器中，原始音频采样率 16\,kHz，使用 STFT 提取 Log-Mel 梅尔频谱，帧移 Hop Size $=160$ 个采样点；得到梅尔频谱特征序列后，经过两层步长 Stride$=2$ 的一维卷积降采样，送入 Transformer。若一段音频总长 12 秒，求最终输入 Transformer 的 Token 序列长度 $L$：}{}
\fourchoicestwo{$1200$}{$600$}{$300$}{$150$}

\qitem{11}{文生图领域存在过多种经典底层生成范式：GAN、VAE、离散自回归 Transformer、扩散 Diffusion 模型。下列关于这些架构的描述，\textbf{错误}的是}{}
\fourchoices{GAN 类文生图（AttnGAN、StackGAN）依靠生成器与判别器对抗训练，单步直接输出图像，无迭代去噪流程；核心痛点是训练易模式崩溃、高分辨率细节生成能力弱}{VAE 类文生图（早期 DALL$\cdot$E 1）基于潜空间概率建模，通过编码器压缩图像、解码器依据文本重建画面；天然存在生成图像偏模糊、纹理表现力不足的缺陷}{离散自回归 Transformer 文生图（Parti、Muse）将图像映射为 VQ-VAE 离散码本，像文本一样逐 Token 串行预测生成；不依赖扩散前向加噪、反向去噪流程}{主流扩散文生图模型里，SD 1.5/XL 的 U-Net 主干、SD3 与 FLUX 的 MMDiT 主干，采用 AdaLN-Zero 完成文本条件的特征调制}

\qitem{12}{在 CLIP 批内对比学习中，设 batch 大小为 $B$，图像嵌入为 $\{\mathbf{I}_i\}_{i=1}^{B}$，文本嵌入为 $\{\mathbf{T}_i\}_{i=1}^{B}$（均已 L2 归一化），相似度 logits 为 $s_{ij}=\mathbf{I}_i^{\top}\mathbf{T}_j/\tau$。下列说法\textbf{正确}的是}{}
\fourchoices{温度 $\tau$ 越小，softmax 分布越平滑，负样本梯度越弱}{对角线元素 $s_{ii}$ 为正样本对，非对角线 $s_{ij}\,(i\neq j)$ 在图像到文本方向充当 in-batch 负样本；$\tau$ 越小分布越尖锐，对困难负样本的区分力越强}{对比矩阵只需计算上三角，因为 $s_{ij}=s_{ji}$ 恒成立}{$B=1$ 时 InfoNCE 仍能有效训练，因为预训练数据足够大}

% ============ 二、填空题 ============
\examsection{二、填空题}{本大题共 6 小题，每小题 5 分，共 30 分。把答案填在题中横线上。}

\qitem{13}{设 ViT 输入图像分辨率为 $H=W=384$，patch 大小 $P=16$，使用 1 个 CLS token。则 patch 数量为 $N_v=$ \blank{2cm}，送入 Transformer 的总序列长度为 \blank{2cm}。}{}

\qitem{14}{在 DDIM 中，设噪声调度满足 $\alpha_t=1-0.01t$（$t=0,1,\ldots,100$），且取 $\sigma_t=0$。已知第 50 步的状态 $x_{50}$ 和模型预测的噪声 $\epsilon$，则计算 $x_{49}$ 时，$x_{50}$ 前面的系数为 \blank{3cm}。}{}

\qitem{15}{在 CLIP 批内对比学习中，batch 大小为 $B$ 时，对于样本 $i$，图像到文本方向的 in-batch 负样本数量为 \blank{2cm}；图像—文本相似度矩阵规模为 \blank{2.5cm}。}{}

\qitem{16}{设视觉 token 数 $N_v=576$，隐藏维度 $d=1024$，注意力头数 $h=16$。单层视觉自注意力中，$QK^{\top}$ 矩阵的显存占用（fp16，每元素 2 字节）约为 \blank{3cm} MB（保留整数）。}{}

\qitem{17}{DDPM 前向过程定义 $\alpha_t=1-\beta_t$，$\bar\alpha_t=\prod_{s=1}^{t}\alpha_s$。则 $q(x_t|x_0)$ 的闭式分布为
\[
q(x_t|x_0)=\mathcal{N}\!\left(\blank{4cm},\;\blank{4cm}\right).
\]
}{}

\qitem{18}{BLIP-2 中 Q-Former 使用 $N_q$ 个可学习 query token 通过交叉注意力从冻结的视觉编码器输出中抽取特征。若 $N_q=32$，视觉编码器输出序列长度为 $N_v=257$（含 CLS），则 Q-Former 输出的视觉 token 数为 \blank{2cm}；相比直接将 $N_v$ 个 token 送入 LLM，LLM 侧序列长度减少比例为 \blank{3cm}（保留两位小数）。}{}

\needspace{1cm}

% ============ 三、解答题 ============
\examsection{三、解答题}{本大题共 6 小题，共 60 分。解答应写出文字说明、证明过程或演算步骤。}

\qitem{19}{（本小题满分 10 分）}{}
\noindent CLIP 通过批内对比学习实现图文对齐。设 batch 大小为 $B$，L2 归一化后的图像嵌入为 $\mathbf{I}_i$，文本嵌入为 $\mathbf{T}_i$，温度系数为 $\tau$，相似度 $s_{ij}=\mathbf{I}_i^{\top}\mathbf{T}_j/\tau$。
\begin{enumerate}[label=(\arabic*),leftmargin=2em,itemsep=0.2em,parsep=0pt,topsep=0pt]
  \item 写出 CLIP 对称 InfoNCE 损失的完整表达式（图像$\to$文本与文本$\to$图像两个方向）；（3 分）
  \item 分析温度系数 $\tau$ 对 softmax 分布与梯度的影响：$\tau$ 增大和减小时分别有何效果？（3 分）
  \item 对图像到文本方向，求 $\mathcal{L}_{I\to T}$ 关于正样本 logit $s_{ii}$ 的梯度 $\partial\mathcal{L}_{I\to T}/\partial s_{ii}$，并解释其直觉含义。（4 分）
\end{enumerate}

\answerlines{6}

\needspace{3cm}
\qitem{20}{（本小题满分 10 分）}{}
\noindent 设 ViT 输入图像 $H\times W$，patch 大小 $P\times P$，隐藏维度 $d$，注意力头数 $h$，每头维度 $d_h=d/h$，编码器层数 $L_{\text{enc}}$。
\begin{enumerate}[label=(\arabic*),leftmargin=2em,itemsep=0.2em,parsep=0pt,topsep=0pt]
  \item 推导 patch 数量 $N_v$ 与总序列长度 $N$（含 CLS token）的表达式；（2 分）
  \item 推导单层 ViT 编码器中 Self-Attention 前向传播的 FLOPs（含 $Q,K,V$ 投影、$QK^{\top}$、$\mathrm{softmax}(QK^{\top})V$、输出投影），用 $N,d,h$ 表示；（4 分）
  \item 若图像分辨率从 $224\times 224$ 提升至 $448\times 448$（$P$ 不变），分析 $N$、单层注意力 FLOPs 及注意力矩阵显存的变化倍数。（4 分）
\end{enumerate}

\answerlines{8}

\needspace{3cm}
\qitem{21}{（本小题满分 10 分）}{}
\noindent 在多模态大模型中，为了同时处理文本、图像和视频，常采用多维旋转位置编码（RoPE）。假设我们有一个 $2\times 2$ 的图像，像素坐标 $(h,w)$ 分别表示行和列，$h,w\in\{0,1\}$。每个像素的特征向量为 4 维，前 2 维用高度编码，后 2 维用宽度编码。1D RoPE 对二维子向量的变换定义为
\[
R_\theta\!\begin{pmatrix}a\\b\end{pmatrix}
=\begin{pmatrix}a\cos\theta-b\sin\theta\\ a\sin\theta+b\cos\theta\end{pmatrix}.
\]
设高度基频 $\omega_h=4\pi$，宽度基频 $\omega_w=2\pi$。
\begin{enumerate}[label=(\arabic*),leftmargin=2em,itemsep=0.2em,parsep=0pt,topsep=0pt]
  \item 已知某像素位于坐标 $(1,0)$，其原始特征向量为 $\mathbf{v}=[1,0,0,1]^{\top}$，求经过 2D RoPE 编码后的向量。（3 分）
  \item 已知一段 2 帧视频，每帧大小为 $2\times 2$，体素坐标 $(t,h,w)$，$t\in\{0,1\}$。现采用 M-RoPE，将特征向量拆分为时间、高度、宽度三组，每组 2 维。基频分别为 $\omega_t=3\pi$，$\omega_h=4\pi$，$\omega_w=2\pi$。某体素位于 $(1,0,1)$，原始特征向量为 $\mathbf{u}=[1,0,0,1,1,0]^{\top}$（依次对应时间、高度、宽度组），求 M-RoPE 编码后的向量。（3 分）
  \item 将坐标 $x$ 视为连续变量，频率 $\omega$ 取遍所有实数，且各频率权重相等。于是位置 $x$ 的编码向量（忽略幅度）可视为函数 $\phi(x)=e^{i\omega x}$ 关于 $\omega$ 的连续族。对于任意两个位置 $x$ 和 $y$，求它们的编码内积 $K(x,y)$。（4 分）
\end{enumerate}

\answerlines{12}

\needspace{6cm}
\qitem{22}{（本小题满分 10 分）}{}
\noindent LLaVA 将预训练视觉编码器（如 CLIP ViT）与 LLM 通过轻量投影器连接，采用两阶段训练策略。
\begin{enumerate}[label=(\arabic*),leftmargin=2em,itemsep=0.2em,parsep=0pt,topsep=0pt]
  \item 画出 LLaVA 前向数据流：从原始图像到 LLM 输出的完整路径，标注视觉编码器、投影器、LLM 各阶段的输入输出维度变化；（4 分）
  \item 对比阶段一（特征对齐预训练）与阶段二（端到端指令微调）的训练目标、冻结策略与数据形式；（3 分）
  \item 若视觉编码器输出 $N_v=576$ 个 token，LLM 最大上下文 4096 token，用户文本平均 200 token。估算视觉 token 占上下文的比例，并说明为何高分辨率场景下需要 AnyRes 或 token 压缩。（3 分）
\end{enumerate}

\answerlines{5}

\needspace{6cm}
\qitem{23}{（本小题满分 10 分）}{}
\noindent BLIP-2 的 Q-Former 与 Flamingo 的 Perceiver Resampler 均通过可学习 query 压缩视觉特征。
\begin{enumerate}[label=(\arabic*),leftmargin=2em,itemsep=0.2em,parsep=0pt,topsep=0pt]
  \item 描述 Q-Former 的结构：可学习 query 如何通过 Cross-Attention 从冻结 ViT 输出 $H_v\in\mathbb{R}^{N_v\times d_v}$ 中提取 $N_q$ 个视觉 token；（4 分）
  \item 对比 Q-Former 与 Flamingo Perceiver Resampler 的异同（输入、注意力类型、输出用途）；（3 分）
  \item 设 $N_v=1024$，$N_q=64$，$d=4096$，LLM 层数 $L=32$。估算将 $N_v$ 个 token 直接送入 LLM 与经 Q-Former 压缩后，LLM 侧 KV Cache 显存（fp16）的节省倍数。（3 分）
\end{enumerate}

\answerlines{5}

\needspace{6cm}
\qitem{24}{（本小题满分 10 分）}{}
\noindent 设 DDPM 前向过程为 $q(x_t|x_{t-1})=\mathcal{N}(\sqrt{1-\beta_t}\,x_{t-1},\,\beta_t I)$，$\alpha_t=1-\beta_t$，$\bar\alpha_t=\prod_{s=1}^{t}\alpha_s$。
\begin{enumerate}[label=(\arabic*),leftmargin=2em,itemsep=0.2em,parsep=0pt,topsep=0pt]
  \item 利用马尔可夫性质与重参数化技巧，推导 $q(x_t|x_0)=\mathcal{N}(\sqrt{\bar\alpha_t}\,x_0,\,(1-\bar\alpha_t)I)$ 的完整过程；（5 分）
  \item 说明为何训练目标可简化为 $\mathcal{L}_{\text{simple}}=\mathbb{E}_{t,x_0,\epsilon}\big[\|\epsilon-\epsilon_\theta(x_t,t)\|^2\big]$，其中 $x_t=\sqrt{\bar\alpha_t}\,x_0+\sqrt{1-\bar\alpha_t}\,\epsilon$，$\epsilon\sim\mathcal{N}(0,I)$。（5 分）
\end{enumerate}

\answerlines{12}

% ============ 答案另起一页 ============
\newpage
\begin{center}
{\zihao{2}\heiti 参考答案与解析}\\[0.3cm]
{\small 命题：大模型多模态课程组 \quad 校对：教研组}
\end{center}
\vspace{0.3cm}
\noindent\rule{\linewidth}{0.6pt}

\subsection*{\heiti 一、选择题答案}

\begin{multicols}{4}
\begin{enumerate}[label=\textbf{\arabic*.},leftmargin=*,itemsep=0.1em,parsep=0pt,topsep=0pt]
  \item A
  \item B
  \item C
  \item B
  \item D
  \item B
  \item B
  \item A
  \item D
  \item C
  \item D
  \item B
\end{enumerate}
\end{multicols}

\vspace{0.3cm}
\noindent\textbf{1. A}\quad 多模态三要素：表示解决"如何编码各模态"；对齐解决"不同模态语义如何对应"；融合解决"多模态信息如何联合用于下游任务"。三者递进且相互关联。B 错误：三者含义不同；C 错误：推理时仍需 cross-attention 等融合；D 错误：VQA 等生成任务同样需要对齐。

\noindent\textbf{2. B}\quad CLIP 采用图像编码器与文本编码器双塔结构，批内 $B$ 个图文对构成 $B\times B$ 相似度矩阵，对角线为正样本、非对角线为负样本，用 InfoNCE 对比损失训练。A 描述的是单塔融合架构；C 是像素级重建而非语义对齐；D 错误：ViT 可通过插值泛化到其他分辨率。

\noindent\textbf{3. C}\quad Patch 数量 $N_v=(224/16)^2=14^2=196$，加 1 个 CLS token 后总序列长度 $N=196+1=197$。

\noindent\textbf{4. B}\quad LLaVA 用线性或 MLP 将 ViT 输出投影到 LLM 词嵌入空间；BLIP-2 用 Q-Former（可学习 query + 交叉注意力）桥接冻结 ViT 与 LLM；Flamingo 在 LLM 层间插入门控交叉注意力，使文本 token 查询视觉特征。

\noindent\textbf{5. D}\quad Late Fusion 虽简单，但无法在中间层捕获细粒度跨模态交互（如"图中红色物体"与"the red object"的 token 级对齐），Intermediate Fusion 在 VQA、Caption 等任务上通常更优。A、B、C 均为正确描述。

\noindent\textbf{6. B}\quad LDM 先用 VAE 编码器将图像压缩为低维 $z$（如 $64\times 64\times 4$），在 $z$ 空间训练 U-Net 去噪，解码时用 VAE 解码器还原图像，计算量远低于像素空间扩散。

\noindent\textbf{7. B}\quad 前向过程是马尔可夫链；当 $t\to T$，$\bar\alpha_t\to 0$，$x_t$ 近似 $\mathcal{N}(0,I)$。反向 $q(x_{t-1}|x_t)$ 未知，需学习 $p_\theta(x_{t-1}|x_t)$；$\beta_t$ 通常采用线性或 cosine 调度，非常数。

\noindent\textbf{8. A}\quad VQ 将连续特征量化为 codebook 索引，argmax 不可微，常用 Straight-Through：前向用离散索引，反向梯度直通编码器输出。B 错误：DALL-E、Parti 等用离散 token + Transformer；C 错误：VQ 损失含 codebook、commitment、重建项；D 错误：码本过大易导致利用率低。

\noindent\textbf{9. D}\quad FLUX 采用 Rectified Flow / Flow Matching 训练目标，区别于经典 DDPM 噪声预测损失；FLUX 为 Transformer 扩散主干（MMDiT 风格），非 SD 系列 U-Net；FLUX 在文字渲染、手部等细节场景口碑普遍优于 SDXL；FLUX.1 系列参数量与激活显存开销通常\textbf{高于} SDXL，本地推理常见需求为 12GB 及以上显存。

\noindent\textbf{10. C}\quad 采样率 $16\,\text{kHz}=16000\,\text{Hz}$，帧移 $160$ 采样点，梅尔帧率 $=16000/160=100$ 帧/秒。12 秒音频对应 $12\times 100=1200$ 帧 Mel 特征。经两层 Stride$=2$ 一维卷积，长度依次减半：$1200\to 600\to 300$，故 $L=300$，选 C。A 为未卷积降采样长度；B 为仅一层 Stride$=2$；D 为三层降采样结果。

\noindent\textbf{11. D}\quad A、B、C 均为正确描述：GAN 对抗单步生成、易模式崩溃；VAE/早期 DALL$\cdot$E 潜空间重建偏模糊；Parti/Muse 等 VQ 离散 token 自回归、不经扩散去噪。D 错误：\textbf{AdaLN-Zero} 是 DiT（Diffusion Transformer）中的文本/条件调制方式；SD 1.5/XL 的 U-Net 主干主要通过\textbf{交叉注意力（Cross-Attention）}注入文本条件，而非 AdaLN-Zero。SD3 的 MMDiT、FLUX 等虽为 Transformer 扩散主干，其条件注入机制也与「全体主流模型均借助 AdaLN-Zero」的表述不符。

\noindent\textbf{12. B}\quad 对角线 $s_{ii}$ 为正样本；$\tau$ 越小 softmax 越尖锐，对困难负样本梯度越大。A 错误：$\tau$ 小则分布更尖锐；C 错误：图像—文本矩阵一般不对称（$s_{ij}\neq s_{ji}$ 当使用不同编码器时）；D 错误：$B=1$ 无负样本，对比学习失效。

\subsection*{\heiti 二、填空题答案}

\noindent\textbf{13.}\quad $N_v=(384/16)^2=24^2=576$\quad ；\quad 总序列长度 $N=576+1=577$。

\noindent\textbf{14.}\quad $\sqrt{\dfrac{\alpha_{49}}{\alpha_{50}}}=\sqrt{\dfrac{0.51}{0.5}}=\sqrt{\dfrac{51}{50}}\approx 1.01$。

\textbf{推导：}DDIM 在 $\sigma_t=0$（确定性采样）时，
\[
x_{t-1}=\sqrt{\alpha_{t-1}}\,\hat x_0+\sqrt{1-\alpha_{t-1}}\,\epsilon,\qquad
\hat x_0=\frac{x_t-\sqrt{1-\alpha_t}\,\epsilon}{\sqrt{\alpha_t}}.
\]
代入并整理 $x_t$ 的系数：
\[
x_{t-1}=\sqrt{\frac{\alpha_{t-1}}{\alpha_t}}\,x_t+\left(\sqrt{1-\alpha_{t-1}}-\sqrt{\frac{\alpha_{t-1}(1-\alpha_t)}{\alpha_t}}\right)\epsilon.
\]
故 $x_{50}$ 的系数为 $\sqrt{\alpha_{49}/\alpha_{50}}$。由 $\alpha_t=1-0.01t$ 得 $\alpha_{50}=0.5$，$\alpha_{49}=0.51$，系数 $=\sqrt{0.51/0.5}=\sqrt{51/50}$。

\noindent\textbf{15.}\quad 负样本数 $=B-1$\quad ；\quad 相似度矩阵规模 $=B\times B$。

\noindent\textbf{16.}\quad $QK^{\top}$ 规模为 $576\times 576$，fp16 显存 $=576^2\times 2=663\,552$ 字节 $\approx 0.66$ MB，保留整数为 $\boxed{1}$ MB。

\noindent\textbf{17.}\quad 均值 $\sqrt{\bar\alpha_t}\,x_0$\quad ；\quad 方差 $(1-\bar\alpha_t)I$（或协方差矩阵 $(1-\bar\alpha_t)I$）。

\noindent\textbf{18.}\quad 输出视觉 token 数 $=32$\quad ；\quad 减少比例 $=1-32/257\approx 0.88$（或 $(257-32)/257\approx 0.88$）。

\subsection*{\heiti 三、解答题答案}

\subsubsection*{19. CLIP 对比学习（10 分）}

\noindent\textbf{(1) 对称 InfoNCE（3 分）}

\[
\mathcal{L}_{I\to T}=-\frac{1}{B}\sum_{i=1}^{B}\log\frac{\exp(s_{ii})}{\sum_{j=1}^{B}\exp(s_{ij})},\quad
\mathcal{L}_{T\to I}=-\frac{1}{B}\sum_{j=1}^{B}\log\frac{\exp(s_{jj})}{\sum_{i=1}^{B}\exp(s_{ij})},
\]
\[
\mathcal{L}_{\text{CLIP}}=\frac{1}{2}(\mathcal{L}_{I\to T}+\mathcal{L}_{T\to I}).
\]

\noindent\textbf{(2) 温度 $\tau$ 的影响（3 分）}

$\tau$ 增大：logits $s_{ij}/\tau$ 变小，softmax 分布更平滑，所有样本（含负样本）获得更均匀的梯度，学习更稳定但区分力下降。$\tau$ 减小：分布更尖锐，模型更关注困难负样本（高相似度但不匹配的对），梯度集中在少数样本，有利于细粒度对齐但可能训练不稳定。CLIP 典型 $\tau\approx 0.07$。

\noindent\textbf{(3) 正样本 logit 梯度（4 分）}

记 $p_{ij}=\exp(s_{ij})/\sum_k\exp(s_{ik})$，则 $\mathcal{L}_{I\to T}=-\frac{1}{B}\sum_i\log p_{ii}$，故
\[
\frac{\partial\mathcal{L}_{I\to T}}{\partial s_{ii}}=-\frac{1}{B}(1-p_{ii}).
\]
直觉：$p_{ii}$ 越接近 1（正样本已很好区分），梯度越接近 0；$p_{ii}$ 小时梯度为负且绝对值大，推动 $s_{ii}$ 增大，拉近正样本对。

\subsubsection*{20. ViT 计算量推导（10 分）}

\noindent\textbf{(1) 序列长度（2 分）}

$N_v=\dfrac{H}{P}\cdot\dfrac{W}{P}$，总序列长度 $N=N_v+1$（含 CLS）。

\noindent\textbf{(2) 单层 Self-Attention FLOPs（4 分）}

$Q,K,V$ 投影：各 $2Nd^2$，共 $6Nd^2$。$S=QK^{\top}$：$2N^2 d$。$O=\mathrm{softmax}(S)V$：$2N^2 d$。输出投影：$2Nd^2$。单层注意力合计 $\approx 8Nd^2+4N^2 d$ FLOPs。

\noindent\textbf{(3) 分辨率翻倍的影响（4 分）}

$224\to 448$：$N_v$ 从 $14^2=196$ 增至 $28^2=784$，约 4 倍。$N$ 从 197 增至 785。注意力 FLOPs 中 $N^2 d$ 项增约 $4^2=16$ 倍，$Nd^2$ 项增 4 倍；显存 $N\times N$ 注意力矩阵增 16 倍。故高分辨率是视觉 token 爆炸的主要来源。

\subsubsection*{21. 2D RoPE / M-RoPE 计算与连续核（10 分）}

\noindent\textbf{(1) 2D RoPE 编码（3 分）}

像素坐标 $(h,w)=(1,0)$，$\mathbf{v}=[v_1,v_2,v_3,v_4]^{\top}=[1,0,0,1]^{\top}$。2D RoPE 对高度、宽度子向量分别施加旋转，旋转角为坐标乘以对应基频：
\[
\theta_h=h\,\omega_h=1\times 4\pi=4\pi,\qquad
\theta_w=w\,\omega_w=0\times 2\pi=0.
\]

\textbf{高度组}（前两维 $[1,0]^{\top}$）：
\[
R_{4\pi}\!\begin{pmatrix}1\\0\end{pmatrix}
=\begin{pmatrix}\cos 4\pi\cdot 1-\sin 4\pi\cdot 0\\ \sin 4\pi\cdot 1+\cos 4\pi\cdot 0\end{pmatrix}
=\begin{pmatrix}1\\0\end{pmatrix}.
\]

\textbf{宽度组}（后两维 $[0,1]^{\top}$）：
\[
R_{0}\!\begin{pmatrix}0\\1\end{pmatrix}
=\begin{pmatrix}\cos 0\cdot 0-\sin 0\cdot 1\\ \sin 0\cdot 0+\cos 0\cdot 1\end{pmatrix}
=\begin{pmatrix}0\\1\end{pmatrix}.
\]

拼接得
\[
\boxed{\mathbf{v}'=[1,\,0,\,0,\,1]^{\top}.}
\]
（本例中 $4\pi$ 与 $0$ 角均为 $2\pi$ 整数倍，旋转退化为恒等变换。）

\noindent\textbf{(2) M-RoPE 编码（3 分）}

体素坐标 $(t,h,w)=(1,0,1)$，$\mathbf{u}=[u_1,u_2,u_3,u_4,u_5,u_6]^{\top}=[1,0,0,1,1,0]^{\top}$。三组旋转角：
\[
\theta_t=t\,\omega_t=3\pi,\qquad
\theta_h=h\,\omega_h=0,\qquad
\theta_w=w\,\omega_w=2\pi.
\]

\textbf{时间组} $[1,0]^{\top}$：$\cos 3\pi=-1$，$\sin 3\pi=0$，
\[
R_{3\pi}\!\begin{pmatrix}1\\0\end{pmatrix}
=\begin{pmatrix}-1\\0\end{pmatrix}.
\]

\textbf{高度组} $[0,1]^{\top}$：$\theta_h=0$，
\[
R_{0}\!\begin{pmatrix}0\\1\end{pmatrix}
=\begin{pmatrix}0\\1\end{pmatrix}.
\]

\textbf{宽度组} $[1,0]^{\top}$：$\cos 2\pi=1$，$\sin 2\pi=0$，
\[
R_{2\pi}\!\begin{pmatrix}1\\0\end{pmatrix}
=\begin{pmatrix}1\\0\end{pmatrix}.
\]

拼接得
\[
\boxed{\mathbf{u}'=[-1,\,0,\,0,\,1,\,1,\,0]^{\top}.}
\]

\noindent\textbf{(3) 连续频率编码内积 $K(x,y)$（4 分）}

\textbf{Step 1：定义内积核。}
\[
K(x,y)=\int_{-\infty}^{\infty}\phi(x,\omega)\,\overline{\phi(y,\omega)}\,\mathrm{d}\omega
=\int_{-\infty}^{\infty}e^{i\omega x}\,e^{-i\omega y}\,\mathrm{d}\omega
=\int_{-\infty}^{\infty}e^{i\omega(x-y)}\,\mathrm{d}\omega.
\]

\textbf{Step 2：傅里叶逆变换。}

在分布（广义函数）意义下：
\[
\int_{-\infty}^{\infty}e^{i\omega(x-y)}\,\mathrm{d}\omega=2\pi\,\delta(x-y).
\]
其中 $\delta(\cdot)$ 为\textbf{狄拉克函数}。

\textbf{Step 3：结论与意义。}
\[
\boxed{K(x,y)=2\pi\,\delta(x-y).}
\]
内积仅依赖\textbf{相对位置 $x-y$}。离散 RoPE 用有限频率 $\theta_k$ 逼近该核，得 $K(x,y)\approx\sum_k\cos((x-y)\theta_k)$，与 (1)(2) 中 $\cos\theta$、$\sin\theta$ 的旋转形式一致。（推导毕）

\end{document}

\subsubsection*{22. LLaVA 架构与训练（10 分）}

\noindent\textbf{(1) 数据流（4 分）}

图像 $\to$ ViT $\to$ $[N_v\times d_v]$ $\to$ 投影器（Linear/MLP）$\to$ $[N_v\times d_{\text{LLM}}]$ $\to$ 与文本 token 嵌入拼接 $\to$ LLM 自回归生成。投影器将视觉特征映射到 LLM 词嵌入空间。

\noindent\textbf{(2) 两阶段对比（3 分）}

阶段一：冻结 ViT 与 LLM，只训练投影器，用图文对（如 CC3M）做 next-token 预测，学习视觉—语言对齐。阶段二：解冻 LLM（或 LoRA），用指令数据（VQA、对话）端到端微调，学习遵循指令与推理。

\noindent\textbf{(3) Token 预算（3 分）}

视觉占比 $=576/(576+200)\approx 74\%$。高分辨率时 $N_v$ 可超 2k，易占满 4k 上下文，需 AnyRes 切图后选择性编码，或 2$\times$2 pooling 合并 token，或 Q-Former 压缩。

\subsubsection*{23. Q-Former 与 Perceiver Resampler（10 分）}

\noindent\textbf{(1) Q-Former 结构（4 分）}

$N_q$ 个可学习 query $\mathbf{Q}\in\mathbb{R}^{N_q\times d}$ 作为 Query，ViT 输出 $H_v$ 作为 Key/Value，经 Cross-Attention：$\text{Attn}(\mathbf{Q},H_v,H_v)$，输出 $N_q$ 个视觉 token。ViT 冻结，只训练 Q-Former 与后续 LLM 接口。

\noindent\textbf{(2) 异同（3 分）}

Q-Former：独立 BERT 式模块，输出接 LLM 投影；Flamingo Perceiver Resampler：在 Flamingo 框架内，用固定数量 latent 查询视觉特征，输出注入 LLM 交叉注意力层。二者均为"少量 query + Cross-Attn"压缩，但集成方式不同。

\noindent\textbf{(3) KV Cache 节省（3 分）}

LLM KV Cache 与序列长度成正比。直接输入：$N_v=1024$；压缩后：$N_q=64$。节省倍数 $=1024/64=16$。32 层 fp16：约 $2\times L\times N\times d\times 2$ 字节，序列减 16 倍则 KV Cache 减 16 倍。


\subsubsection*{24. 扩散模型推导（10 分）}

\noindent\textbf{(1) $q(x_t|x_0)$ 闭式解推导（5 分）}

\textbf{Step 1：写出单步重参数化形式。}

由题设前向马尔可夫过程
\[
q(x_t|x_{t-1})=\mathcal{N}\!\left(\sqrt{1-\beta_t}\,x_{t-1},\,\beta_t I\right),
\quad \alpha_t=1-\beta_t,
\]
可重参数化为
\[
x_t=\sqrt{\alpha_t}\,x_{t-1}+\sqrt{1-\alpha_t}\,\epsilon_{t-1},\qquad \epsilon_{t-1}\sim\mathcal{N}(0,I).
\]

\textbf{Step 2：递推至 $x_1$。}

\[
x_1=\sqrt{\alpha_1}\,x_0+\sqrt{1-\alpha_1}\,\epsilon_0.
\]

\textbf{Step 3：递推至 $x_2$ 并观察系数规律。}

\[
\begin{aligned}
x_2&=\sqrt{\alpha_2}\,x_1+\sqrt{1-\alpha_2}\,\epsilon_1\\
&=\sqrt{\alpha_2}\!\left(\sqrt{\alpha_1}\,x_0+\sqrt{1-\alpha_1}\,\epsilon_0\right)+\sqrt{1-\alpha_2}\,\epsilon_1\\
&=\sqrt{\alpha_1\alpha_2}\,x_0+\underbrace{\sqrt{\alpha_2(1-\alpha_1)}\,\epsilon_0+\sqrt{1-\alpha_2}\,\epsilon_1}_{\text{两个独立高斯的线性组合}}.
\end{aligned}
\]
因 $\epsilon_0,\epsilon_1$ 独立标准高斯，其线性组合仍为零均值高斯，方差为
\[
\alpha_2(1-\alpha_1)+(1-\alpha_2)=1-\alpha_1\alpha_2.
\]
故可写为 $x_2=\sqrt{\bar\alpha_2}\,x_0+\sqrt{1-\bar\alpha_2}\,\bar\epsilon_1$，其中 $\bar\alpha_2=\alpha_1\alpha_2$，$\bar\epsilon_1\sim\mathcal{N}(0,I)$。

\textbf{Step 4：数学归纳法推广到一般 $t$。}

设 $\bar\alpha_t=\prod_{s=1}^{t}\alpha_s$，且已有
\[
x_{t-1}=\sqrt{\bar\alpha_{t-1}}\,x_0+\sqrt{1-\bar\alpha_{t-1}}\,\bar\epsilon_{t-2}.
\]
则
\[
\begin{aligned}
x_t&=\sqrt{\alpha_t}\,x_{t-1}+\sqrt{1-\alpha_t}\,\epsilon_{t-1}\\
&=\sqrt{\alpha_t\bar\alpha_{t-1}}\,x_0+\sqrt{\alpha_t(1-\bar\alpha_{t-1})}\,\bar\epsilon_{t-2}+\sqrt{1-\alpha_t}\,\epsilon_{t-1}.
\end{aligned}
\]
注意到 $\alpha_t\bar\alpha_{t-1}=\bar\alpha_t$。噪声项方差为
\[
\alpha_t(1-\bar\alpha_{t-1})+(1-\alpha_t)=1-\alpha_t\bar\alpha_{t-1}=1-\bar\alpha_t.
\]
因此
\[
x_t=\sqrt{\bar\alpha_t}\,x_0+\sqrt{1-\bar\alpha_t}\,\epsilon,\qquad \epsilon\sim\mathcal{N}(0,I).
\]
即任意 $t$ 步后，条件分布为
\[
\boxed{q(x_t|x_0)=\mathcal{N}\!\left(\sqrt{\bar\alpha_t}\,x_0,\,(1-\bar\alpha_t)I\right).}
\]
当 $t\to T$ 且 $\bar\alpha_T\to 0$ 时，$x_T\approx\mathcal{N}(0,I)$，实现向纯噪声的渐进破坏。（推导毕）

\noindent\textbf{(2) 训练目标简化为噪声预测（5 分）}

\textbf{Step 1：变分下界的原始形式。}

DDPM 训练目标是最大化数据对数似然 $\log p_\theta(x_0)$，等价于最小化变分下界（ELBO）。对每一步去噪分布 $p_\theta(x_{t-1}|x_t)$ 与真实后验 $q(x_{t-1}|x_t,x_0)$ 的 KL 散度可写为（省略常数项）：
\[
\mathcal{L}_{\text{VLB}}=\mathbb{E}_q\!\left[\sum_{t=1}^{T} D_{\mathrm{KL}}\!\big(q(x_{t-1}|x_t,x_0)\,\|\,p_\theta(x_{t-1}|x_t)\big)\right]+\text{const}.
\]
其中真实后验 $q(x_{t-1}|x_t,x_0)$ 在已知 $x_0$ 时有闭式高斯解，其均值是 $x_0$ 与 $x_t$ 的线性组合。

\textbf{Step 2：由 (1) 建立 $x_0$ 与 $\epsilon$ 的一一对应。}

对采样 $t$、数据 $x_0$、噪声 $\epsilon\sim\mathcal{N}(0,I)$，定义
\[
x_t=\sqrt{\bar\alpha_t}\,x_0+\sqrt{1-\bar\alpha_t}\,\epsilon.
\]
由于 $\bar\alpha_t>0$，可反解
\[
x_0=\frac{1}{\sqrt{\bar\alpha_t}}\!\left(x_t-\sqrt{1-\bar\alpha_t}\,\epsilon\right).
\]
故给定 $(x_t,t)$ 时，预测 $x_0$、预测 $\epsilon$、预测 $x_{t-1}$ 三者\textbf{参数化等价}，可互相线性变换。

\textbf{Step 3：将 KL 项改写为均方误差。}

Ho 等（DDPM）证明：在参数化
\[
p_\theta(x_{t-1}|x_t)=\mathcal{N}\!\big(\mu_\theta(x_t,t),\,\sigma_t^2 I\big)
\]
且选取与 $q$ 匹配的方差 $\sigma_t^2$ 时，第 $t$ 步 KL 项（忽略权重系数）等价于
\[
\left\|\tilde\mu_t(x_t,x_0)-\mu_\theta(x_t,t)\right\|^2,
\]
其中 $\tilde\mu_t$ 为 $q(x_{t-1}|x_t,x_0)$ 的均值。将 $\tilde\mu_t$ 用 $\epsilon$ 表示后，可进一步化为
\[
\left\|\epsilon-\epsilon_\theta(x_t,t)\right\|^2.
\]
即网络直接预测注入噪声 $\epsilon$，与预测均值 $\mu_\theta$ 或预测 $x_0$ 在数学上等价。

\textbf{Step 4：得到 Simple Loss 及其优点。}

对 $t$ 均匀采样，最终常用的简化训练目标为
\[
\boxed{\mathcal{L}_{\text{simple}}=\mathbb{E}_{t\sim\mathcal{U}\{1,\ldots,T\},\,x_0,\,\epsilon}\!\left[\left\|\epsilon-\epsilon_\theta(x_t,t)\right\|^2\right].}
\]
相比完整 VLB：（i）\textbf{实现简单}，只需对加噪样本回归噪声；（ii）\textbf{训练稳定}，各 $t$ 步梯度尺度更均衡；（iii）\textbf{效果相当}，实践中与加权 VLB 生成质量接近。Stable Diffusion 等 latent 扩散模型在隐空间 $z$ 上沿用同一套 $\epsilon$-prediction 范式。（推导毕）